\chapter{Low-Level Geometry and DoF Access}

\section{Motivation}

Once the higher-level ontological split between \texttt{Vertex}, \texttt{Node},
and \texttt{ElementGeometry} was in place, the next question was lower-level:
how much indirection remains in the hot path, and where does the code still pay
avoidable access cost?

The target of this pass was not a redesign of algorithms.  It was a careful
data-access refinement guided by three HPC-oriented principles:
\begin{enumerate}[nosep]
  \item keep the public concepts explicit and semantically correct;
  \item reduce repeated pointer chasing in frequently evaluated geometry code;
  \item preserve value semantics and compile-time structure.
\end{enumerate}

\section{\texttt{Point<Dim>} as a contiguous data source}

\subsection{Previous state}

\texttt{Point<Dim>} already stored coordinates in a
\texttt{std::array<double, Dim>}.  However, the main accessor returned that
array by value.  For a small fixed-size object this is not catastrophic, but it
encourages copies and obscures the fact that the coordinates are already stored
contiguously in memory.

\subsection{Implemented change}

\texttt{Point<Dim>} now exposes:
\begin{itemize}[nosep]
  \item \texttt{coord()} returning a \emph{const reference} to the internal
        array for lvalues,
  \item \texttt{coord\_ref()} as the explicit aliasing form, and
  \item \texttt{data()} returning a raw pointer to the contiguous coordinate
        storage.
\end{itemize}

This is a small but important clarification.  It makes \texttt{Point} suitable
not only as a semantic geometry entity but also as a low-level storage source
for element kernels.

\subsection{Why this matters}

The geometry kernels repeatedly evaluate expressions of the form
\[
  x_i = \sum_a N_a(\bxi)\, x_{a,i},
  \qquad
  J_{ij} = \sum_a x_{a,i}\, \frac{\partial N_a}{\partial \xi_j}.
\]

The coordinates \(x_{a,i}\) are read many times.  Making the contiguous storage
explicit enables later code to cache direct pointers to those values instead of
repeatedly calling higher-level accessors.

\section{\texttt{Node<Dim, Storage>} and direct DoF access}

\subsection{Added direct pointer interface}

\texttt{Node} already exposed \texttt{dof\_index()} as a span, which is a good
high-level interface.  This pass adds \texttt{dof\_data()} for the cases where
the caller already knows the number of DOFs and only needs the raw contiguous
pointer for PETSc interaction.

This is now used in selected PETSc-heavy paths such as boundary-condition and
model bookkeeping code.

\subsection{Array constructors}

\texttt{Node} also gained constructors from
\texttt{std::array<double, Dim>}.  This reduces the need for
\texttt{std::apply}-style reconstruction when a node is materialised from a
geometry object such as \texttt{Vertex<Dim>}.  In practice, this simplifies the
\texttt{Domain} transition code and keeps template instantiation noise lower.

\section{Cached coordinate pointers inside geometric elements}

\subsection{Problem}

After the previous refactor, each concrete geometry family already stored a
geometry-only binding to \texttt{Point/Vertex}.  That fixed the architectural
coupling, but the implementation still fetched coordinates through chains like
\[
  \texttt{points}[a] \rightarrow \texttt{coord()} \rightarrow \texttt{array}[i].
\]

This is acceptable, but it is still more indirect than necessary inside the
innermost geometry loops.

\subsection{Implemented solution}

The concrete families
\begin{itemize}[nosep]
  \item \texttt{LagrangeElement},
  \item \texttt{SimplexElement}, and
  \item \texttt{SerendipityElement}
\end{itemize}
now cache, at bind time, direct pointers to the coordinate storage of each
bound point:
\[
  \texttt{coord\_data\_[a]} = \texttt{point\_a.data()}.
\]

The geometric hot path then uses:
\begin{itemize}[nosep]
  \item \texttt{coord\_data\_[a][i]} in \texttt{coord\_array()},
  \item \texttt{coord\_data\_[a][i]} in \texttt{evaluate\_jacobian()}.
\end{itemize}

This removes a repeated layer of accessor traffic while preserving the semantic
\texttt{Point}/\texttt{Vertex} binding.

\subsection{Why this is safe}

The cached coordinate pointers are non-owning views into the bound geometry
objects.  They are valid under the same condition that already governed the
older point/node pointer bindings:
\begin{center}
the lifetime and address stability of the bound geometry objects must exceed the
lifetime of the bound element geometry.
\end{center}

This is not a new invariant; it is simply made more explicit.

\section{Residual geometry-via-node usage in beam elements}

\subsection{Observation}

Even after the previous phases, \texttt{BeamElement::compute\_frame()} still
read coordinates from \texttt{geometry\_->node\_p(0)} and
\texttt{geometry\_->node\_p(1)}.  That meant the beam frame computation still
depended on analysis nodes even though it is purely geometric.

\subsection{Implemented change}

The beam frame is now built from:
\[
  \texttt{geometry\_->point\_p(0).coord\_ref()},
  \qquad
  \texttt{geometry\_->point\_p(1).coord\_ref()}.
\]

Therefore a beam element can compute its frame from geometry-only bindings
before nodal DOFs are present.

\section{Validation and traceability}

This pass added or extended tests in three places:
\begin{enumerate}[nosep]
  \item \texttt{test\_point.cpp}: verifies aliasing and contiguous pointer
        access for \texttt{Point};
  \item \texttt{test\_node\_dof.cpp}: verifies array construction and
        \texttt{dof\_data()} consistency for \texttt{Node};
  \item \texttt{test\_beam\_element.cpp}: verifies that beam frame construction
        works from geometric vertices alone, without materialising analysis
        nodes.
\end{enumerate}

These tests complement the earlier \texttt{domain\_vertex\_refactor} checks and
make the refactor traceable from ontology down to low-level data access.

\section{Interpretation}

This chapter is intentionally modest in scope.  It does not claim a
microbenchmark-proven global speedup.  What it does provide is a stricter and
cleaner data path:
\begin{itemize}[nosep]
  \item geometry entities expose their contiguous storage explicitly;
  \item analysis entities expose their contiguous DOF storage explicitly;
  \item geometric elements cache the minimal data needed for repeated
        interpolation and Jacobian evaluation;
  \item structural elements use geometry where geometry is enough.
\end{itemize}

That is the right kind of optimisation for this codebase: local, transparent,
semantically aligned, and compatible with the broader compile-time design.
